\documentclass{article}
\usepackage{amsmath,amssymb}
\usepackage[margin=1in]{geometry}

\title{Boundary-Squeeze / Resource-Attractor Principle\\
\large CANDIDATE architecture, not a theorem}
\author{Note on \texttt{j15marti/pnp}, branch \texttt{main}.\\
Status of the whole principle: CANDIDATE. Nothing here is claimed as proved.}
\date{2026-10-04, about 12:07~PM PT}

\begin{document}
\maketitle

\noindent
P versus NP remains open.
This note records a proposal of that date and time under the name
Boundary-Squeeze / Resource-Attractor Principle.
The status of the whole principle is a \textbf{CANDIDATE architecture, not a theorem}.
It does not solve P versus NP, it does not assert \(\mathrm{P}=\mathrm{NP}\) or \(\mathrm{P}\neq\mathrm{NP}\),
and it does not exhibit an algorithm, a matrix, or a numerical exponent that would decide the question.
A 12:06~PM PT PC/RAM follow-up may still be in progress; this note does not use it and does not depend on it.

\section{Coordinate remark}
Let \(T\) be a time bound with \(T(n)>1\) for large \(n\), and define the exponent coordinate
\[
e_T(n)=\frac{\log T(n)}{\log n},
\]
where \(\log\) is any fixed base greater than \(1\).
The choice of base cancels in the polynomial case and does not change which limits are infinite.

If \(T(n)=n^k\) for a finite constant \(k\), then \(e_T(n)=k\) for every \(n>1\).
The exponent stays finite.
Two superpolynomial examples do not:
\[
T(n)=n^{\log\log n}
\quad\text{gives}\quad
e_T(n)=\log\log n\to\infty,
\]
and
\[
T(n)=2^{\sqrt{n}}
\quad\text{gives}\quad
e_T(n)=\sqrt{n}\cdot\frac{\log 2}{\log n}\to\infty.
\]
There is no largest polynomial.
By definition,
\[
\mathrm{P}=\bigcup_{k<\infty} O(n^k),
\]
a union over finite exponents, not a single \(O(n^K)\) for a maximum \(K\).

Compactify the exponent by
\[
z=\frac{1}{1+e_T}.
\]
Whenever \(e_T(n)\to\infty\), one has \(z(n)\to 0\).
That is where the two superpolynomial examples land.
A polynomial of exponent \(k\) sits at the positive point \(z=1/(1+k)\).
The set of those points, as \(k\) runs through the finite exponents, accumulates at \(0\).
So the closure of the polynomial points and the closure of the superpolynomial limit meet at \(z=0\).
In this coordinate, merely taking limits cannot separate \(\mathrm{P}\) from \(\mathrm{NP}\):
the point that superpolynomial rates approach is already an accumulation point of polynomial rates.
This is a remark about this compactification.
It is not a proof that no other argument can separate the classes.

\section{Dynamical squeeze, conditional on \(\mathrm{P}=\mathrm{NP}\)}
\textbf{Hypothesis, not a claim:} assume \(\mathrm{P}=\mathrm{NP}\).
Then some exact algorithm for SAT runs in polynomial time, so it has some finite exponent \(K\).
That \(K\) is unknown.
Nothing below constructs it.

Consider a \emph{hypothetical} uniform map sending an exact SAT algorithm of exponent \(K\) to an exact SAT algorithm of exponent \(K'\) with
\[
K'\le a+bK, \qquad 0<b<1,
\]
the same constants \(a\) and \(b\) at every iteration, and every construction cost charged into the new exponent.
If that inequality may be reapplied, the iterates satisfy
\[
K_t
\le a\frac{1-b^t}{1-b}+b^t K
\xrightarrow[t\to\infty]{}
\frac{a}{1-b}.
\]
The limit \(a/(1-b)\) does not depend on the starting exponent \(K\).

Suppose, in addition, that some unconditional theorem forbids exact polynomial SAT exponents strictly below a number \(L\), and that the fixed point lies past that wall:
\[
\frac{a}{1-b}<L.
\]
Then a large enough iterate is forced below \(L\), into the forbidden region.
That would contradict the unconditional theorem, and hence would contradict the assumption \(\mathrm{P}=\mathrm{NP}\), \emph{if the map existed}.

The map is not proved.
No constants \(a\) and \(b\) are known, no uniform construction is known, and no charging argument is known.
\textbf{Mark this map OPEN.}

\section{Vector form}
Let a resource vector gather four exponents,
\[
r=(t,s,q,a),
\]
for time, space, query, and acquisition.
These are names for coordinates in the proposal.
They are not measured values.
The hypothetical update is an affine inequality
\[
r'\le Mr+b
\]
with spectral radius \(\rho(M)<1\).
If \(M\) is a nonnegative matrix, the comparison recurrence
\(u_{t+1}=Mu_t+b\) is an upper bound on the iterates, and
\[
u_t=M^t r_0+\sum_{i=0}^{t-1}M^i b
\xrightarrow[t\to\infty]{}
(I-M)^{-1}b,
\]
since \(\rho(M)<1\) makes the Neumann series converge.
If that fixed point lies in an unconditional forbidden region \(F\), and the inequality really forces the resource vector into \(F\), that would be a contradiction.

This is not a theorem.
The matrix \(M\) is not constructed.
The offset \(b\) is not constructed.
Nonnegativity of \(M\), which the comparison uses, is not known; if some entries are negative, componentwise comparison can fail, and that case is not claimed.
\textbf{Mark the vector contraction OPEN.}

\section{SAME\_REALIZATION}
Two different representations that approach incompatible limits are not a contradiction.
A limit of one encoding and a limit of another encoding are not the same object unless a theorem says they are the same realization.
The structural coordinates
\[
\frac{\mathrm{rank}\,\Phi}{m},\qquad
\frac{\rho(E_H)}{m},\qquad
\frac{\varepsilon}{m},\qquad
\frac{\tau}{m}
\]
are \textbf{CANDIDATE only}.
This note does not define a proved meaning for \(\Phi\), \(E_H\), \(\varepsilon\), or \(\tau\), and it does not treat a numerical value of any of them as evidence.

\section{The missing theorem}
\begin{center}
\fbox{\parbox{0.92\linewidth}{
\textbf{OPEN.}
For every polynomial-time exact SAT algorithm \(A\), there is a uniformly constructed exact SAT algorithm \(A'\) whose resource vector is contracted by some matrix \(M\) with spectral radius \(\rho(M)<1\), with all construction costs charged to the new vector, with the transformation repeatable, and with the fixed point \((I-M)^{-1}b\) inside a region that an unconditional theorem has already forbidden.
No such \(M\), no such uniform construction, and no such placement of the fixed point are given in this note.
}}
\end{center}

Until that statement is a theorem, the principle does not separate \(\mathrm{P}\) from \(\mathrm{NP}\).
Conditional reasoning from an unproved map is not a proof.

\section{Literature opened for this note}
Only sources actually retrieved on this pass are treated as read.
Each item states what the opened text proves, then whether it is the missing contraction.
It is not.
No paper below is a proof of the Boundary-Squeeze principle.

\subsection{Alternation trading and \(2\cos(\pi/7)\)}
Opened: Ryan Williams, ``Alternation-Trading Proofs, Linear Programming, and Lower Bounds (Extended Abstract),'' arXiv:1001.0746, HTML at \texttt{https://ar5iv.labs.arxiv.org/html/1001.0746}.

What it proves, as stated there.
The best cited time lower bound for SAT on \(n^{o(1)}\)-space algorithms is \(n^{2\cos(\pi/7)-o(1)}\ge n^{1.801}\), from the paper's reference [Wil08]; the decimal is that paper's comparison, not a calculation done here.
The same text's Theorem~3.9 says that if time \(t(n)\) and space \(s(n)\) are both bounded by polynomials, then any SAT algorithm satisfies \(t\cdot s=\Omega(n^{2\cos(\pi/7)-\varepsilon})\) for every \(\varepsilon>0\).
The \(n^{2\cos(\pi/7)}\) bound is tied to an inductive family of proof annotations: with \(A=1,0,1,0,\ldots,1,0,0\), the sequence is \([A]\), \([1,A,A]\), \([1,1,A,A,A]\), and so on.
Speedup in that calculus simulates small-space deterministic time by an alternating machine with a smaller time exponent; slowdown, using an assumed SAT algorithm, removes an alternation and multiplies an exponent.
Theorem~3.8 of that text says there is no alternation-trading proof that \(\mathrm{NTIME}[n]\nsubseteq\mathrm{DTS}[n^2]\).

Whether it is the missing contraction.
No.
The exponent \(2\cos(\pi/7)\) is the limit of a recurrence over \emph{proof annotations} in a space-restricted trading calculus, not the fixed point of a map applied to an unknown polynomial exponent of an arbitrary exact SAT algorithm.
The pure time form is a bound under subpolynomial space, \(n^{o(1)}\), not under unrestricted polynomial space.
The product form forbids pairs of polynomial time and space exponents whose sum is too small.
It does not forbid every large finite time exponent, and it does not build a new SAT algorithm from an old one.
It is a possible shape of forbidden region for a future vector contraction that also controls space.
It is not that contraction.

Opened: Sam Buss and Ryan Williams, ECCC TR11-031, abstract at \texttt{https://eccc.weizmann.ac.il/report/2011/031/}.
The abstract says the paper characterizes alternation-trading proofs that satisfiability is not in \(\mathrm{DTISP}(n^c,n^\varepsilon)\) for various \(c<2\) and \(\varepsilon<1\), and proves Williams' conjecture that \(c=2\cos(\pi/7)\) is optimal for the \(\varepsilon=0\) case with currently known methods.
The full paper was not opened.
This is a limit on a proof system for space-restricted lower bounds.
It is not the missing contraction.

\subsection{Fortnow and Fortnow--Lipton--van Melkebeek--Viglas}
Opened: the author page \texttt{https://pages.cs.wisc.edu/~dieter/Research/sat-ts.html}, which states the abstract of L.~Fortnow, R.~Lipton, D.~van Melkebeek, and A.~Viglas, \emph{Journal of the ACM} 52:835--865, 2005.

What that abstract proves.
For any constant \(c\) less than the golden ratio there is a positive constant \(d\) such that no deterministic random-access Turing machine solves satisfiability in time \(n^c\) and space \(n^d\), and \(d\) approaches \(1\) as \(c\) does.
On conondeterministic machines the same holds for any constant \(c\) less than \(\sqrt{2}\).
The bounds apply to nondeterministic linear time and to the stated class of languages solvable in nondeterministic linear time and space \(n^{1/c}\).
The abstract says the proofs use indirect diagonalization.
This page does not give a decimal for the golden ratio; none is invented here.

Whether it is a forbidden region a contraction could aim at.
It is a forbidden region only together with a space restriction: time exponent \(c\) is ruled out only for some space exponent \(d>0\) depending on \(c\), with \(d\) approaching \(1\) only as \(c\) approaches the golden ratio from below.
An exact SAT algorithm whose space exponent stays large is not placed in that region by the abstract.
A contraction of the time exponent alone, with space left free, need not enter it.
The conondeterministic statement has the same shape, with \(\sqrt{2}\) in place of the golden ratio, and is still space-restricted.
Neither statement maps SAT algorithms to SAT algorithms.
Not the missing contraction.

Fortnow's own paper, ``Time--Space Tradeoffs for Satisfiability,'' was not opened on this pass.
The PDF at \texttt{https://lance.fortnow.com/papers/files/npvsnl.pdf} returned a gateway timeout, the ScienceDirect page timed out, and the IEEE page for the conference version returned a robot check with no abstract.
\textbf{Not verified this pass.}
The opened Williams extended abstract cites Fortnow (1997) and Fortnow--van Melkebeek (2000) and, in Williams' words, says an \(n^\phi\) deterministic lower bound is provable, with \(\phi\) the golden ratio.
That is a citation inside Williams, not a verification of Fortnow's text, and it is not used here as a statement of Fortnow's theorem.

\subsection{Speedup, downward self-reducibility, instance compression}
No source opened on this pass proves a theorem that takes an arbitrary polynomial exact SAT algorithm of exponent \(K\) and returns an exact SAT algorithm of exponent at most \(a+bK\) with \(0<b<1\), costs charged, uniformly, repeatably, and with fixed point past a proved frontier.
That map is still missing.

Opened: Jayalal Sarma, lecture notes, IIT Madras CS6840, 28~January~2012, \texttt{https://www.cse.iitm.ac.in/~jayalal/teaching/CS6840/2012/lecture13.pdf}.
The notes state that satisfiability of a formula on \(n\) variables reduces to satisfiability of two formulas on \(n-1\) variables, and that this yields a polynomial-time construction of a satisfying assignment from a correct decision subroutine, using one pass over the \(n\) variables.
That is search-to-decision.
It is not the missing contraction.
A linear number of calls to a decision procedure of exponent \(K\) produces a search procedure whose exponent is larger by \(1\), not a decision procedure whose exponent is contracted by a factor \(b<1\).

Opened: Danny Harnik and Moni Naor, ECCC TR06-022, abstract at \texttt{https://eccc.weizmann.ac.il/report/2006/022/}, and the HTML abstract at \texttt{http://www.wisdom.weizmann.ac.il/~naor/PAPERS/compressibility.html}.
They define a compressibility question: SAT would be compressible if a formula with \(m\) clauses on \(n\) variables could be replaced by an equisatisfiable formula of size at most \(p(n,\log m)\) for some polynomial \(p\).
The opened abstracts give cryptographic consequences of such compression and limitations on witness-retrievable compression.
They do not prove that SAT is compressible, and they do not produce a slower polynomial time exponent from a faster one.
Not the missing contraction.

Opened, abstract only: Ryan Williams, ECCC TR23-082, ``Self-Improvement for Circuit-Analysis Problems,'' \texttt{https://eccc.weizmann.ac.il/report/2023/082/}.
The abstract states a self-improving theorem for Circuit-SAT, \#Circuit-SAT, and a fully quantified version: a faster algorithm on large circuits implies a speedup on smaller circuits, including non-uniform and randomized models.
It states that certain fine-grained improvements on the exponents of polynomial-time versions of Circuit-SAT would imply subexponential-time Circuit-SAT on \(2^{o(n)}\)-size circuits, refuting the Exponential Time Hypothesis, and it gives one concrete illustration of that shape, an algorithm running in \(1000000^k+n^{1+\varepsilon}\) time for every \(\varepsilon>0\) implying time \((1+\varepsilon)^k+n^{1+\varepsilon}\).
The full paper was not opened.
This bootstraps circuit-analysis algorithms across sizes.
It is not a uniform, cost-charged contraction of every polynomial exact SAT algorithm into a proved polynomial-time forbidden region.
Not the missing contraction.

Not verified this pass, beyond the searches that pointed at them and were not fetched:
the full Fortnow JCSS paper already noted above;
LV-reductions between constraint languages in the fine-grained exponential setting;
and any other claimed ``SAT speedup'' that would contract a polynomial exponent.
Search hits are not treated as opened sources.

\section{Conclusion}
P versus NP remains open.
The compactification explains why a bare limit of \(e_T\) cannot be a separation.
The squeeze, both scalar and vector, is an OPEN shape of argument.
It becomes an argument only if the missing contraction is proved and its fixed point is shown to land in a region an unconditional theorem already forbids, on one and the same realization, with every cost charged.
No such map was found in the sources opened here.
There is no known SAT-specific contraction, among those sources, whose fixed point lies past a proved frontier.
That map is still missing.
\end{document}
